\documentclass{article}
\usepackage[utf8]{inputenc}
\usepackage[T1]{fontenc}
\usepackage[margin=2cm]{geometry}
\usepackage{hyperref}
\sloppy
\begin{document}
\section*{Historical source}
Encyclopaedia Britannica scan, source row 1828. Parse 3.5 extracted the text. Super 3.5 VL corrected it from an enlarged crop of each block. It can still contain recognition errors.
Page headers, footers, and page numbers are removed.

\subsection*{BEAUTY (E0165)}
\textit{Continues an entry from an earlier page.}

duced? and, Secondly. What is the nature of the connexion by which we suppose that the objects we call beautiful are enabled to suggest these affections? With regard to the first of these points, it fortunately is not necessary either to enter into any tedious details, or to have recourse to any nice distinctions. All sensations that are not absolutely indifferent, and are, at the same time, either agreeable when experienced by ourselves, or attractive when contemplated in others, may form the foundation of the emotions of sublimity or beauty. The love of sensation seems to be the ruling appetite of human nature; and many sensations, in which the painful seems to bear no little share, are consequently sought for with avidity, and recollected with interest, even in our own persons. In the persons of others, emotions still more painful are contemplated with eagerness and delight; and therefore we must not be surprised to find, that many of the pleasing sensations of beauty or sublimity resolve themselves ultimately into recollections of feelings that may appear to have a very opposite character. The sum of the whole is, that every feeling which it is agreeable to experience, to recall, or to witness, may become the source of beauty in external objects, when it is so connected with them as that their appearance reminds us of that feeling. Now, in real life, and from daily experience and observation, we know that it is agreeable, in the first place, to recollect our own pleasurable sensations, or to be enabled to form a lively conception of the pleasures of other men, or even of sentient beings of any description. We know, likewise, from the same sure authority, that there is a certain delight in the remembrance of our past, or the conception of our future emotions, even though attended with great pain, provided they be not forced too rudely on the mind, and be softened by the accompaniment of any milder feeling. And, finally, we know, in the same manner, that the spectacle or conception of the emotions of others, even when in a high degree painful, is extremely interesting and attractive, and draws us away, not only from the consideration of indifferent objects, but even from the pursuit of light or frivolous enjoyments. All these are plain and familiar facts, of the existence of which, however they may be explained, no one can entertain the slightest doubt, and into which, therefore, we shall have made no inconsiderable progress, if we can resolve the more mysterious fact, of the emotions we receive from the contemplation of sublimity or beauty.

Our proposition then is, that these emotions are not original emotions, nor produced directly by any qualities in the objects which excite them; but are reflections or images of the more radical and familiar emotions to which we have already alluded, and are occasioned, not by any inherent virtue in the objects before us, but by the accidents, if we may so express ourselves, by which these may have been enabled to suggest or recall to us our own past sensations or sympathies. We might almost venture, indeed, to lay it down as an axiom, that, except in the plain and palpable case of bodily pain or pleasure, we can never be interested in any thing but the fortunes of sentient beings; and that every thing partaking of the nature of mental emotion must have for its object the feelings, past, present, or possible, of something capable of sensation. Independently, therefore, of all evidence, and without the help of any explanation, we should have been apt to conclude that the emotions of beauty and sublimity must have for their objects the sufferings or enjoyments of sentient beings; and to reject as intrinsically absurd and incredible the supposition, that material objects, which obviously do neither hurt nor delight the body, should yet excite, by their mere physical qualities, the very powerful

emotions which are sometimes excited by the spectacle Beauty.
of beauty.

Of the feelings, by their connection with which external objects become beautiful, we do not think it necessary to speak more minutely; and therefore it only remains, under this preliminary view of the subject, to explain the nature of that connection by which we conceive this effect to be produced. Here, also, there is but little need for minuteness, or fulness of enumeration. Almost every tie by which two objects can be bound together in the imagination, in such a manner as that the presentment of the one shall recall the memory of the other, or, in other words, almost every possible relation which can subsist between such objects, may serve to connect the things we call sublime or beautiful, with feelings that are interesting or delightful. It may be useful, however, to class these bonds of association between mind and matter in a rude and general way.

It appears to us, then, that objects are sublime or beautiful, first, When they are the natural signs and perpetual concomitants of pleasurable sensations, or, at any rate, of some lively feeling or emotion in ourselves or in some other sentient beings; or, secondly, When they are the arbitrary or accidental concomitants of such feelings; or, thirdly, When they bear some analogy or fanciful resemblance to things with which these emotions are necessarily connected. In endeavouring to illustrate the nature of these several relations, we shall be led to lay before our readers some proofs that appear to us satisfactory of the truth of the general theory.

The most obvious and the strongest association that can be established between inward feelings and external objects is, where the object is necessarily and universally connected with the feeling by the law of nature, so that it is always presented to the senses when the feeling is impressed upon the mind, as the sight or sound of laughter, with the feeling of gaiety, of weeping, with distress, of the sound of thunder, with ideas of danger and power. Let us dwell for a moment on the last instance. Nothing, perhaps, in the whole range, is more strikingly and universally sublime than the sound we have just mentioned; yet it seems obvious that the sense of sublimity is produced, not by any quality that is perceived by the ear, but altogether by the impression of power and of danger that is necessarily made upon the mind whenever that sound is heard. That it is not produced by any peculiarity in the sound itself, is certain, from the mistakes that are frequently made with regard to it. The noise of a cart rattling over the stones is often mistaken for thunder; and as long as the mistake lasts, this very vulgar and insignificant noise is actually felt to be prodigiously sublime. It is so felt, however, it is perfectly plain, merely because it is then associated with ideas of prodigious power and undefined danger; and the sublimity is destroyed the moment the association is dissolved, though the sound itself, and its effect on the organ, continue exactly the same. This, therefore, is an instance in which sublimity is distinctly proved to consist, not in any physical quality of the object to which it is ascribed, but in its necessary connection with that vast and uncontrolled Power which is the natural object of awe and veneration.

We may now take an example a little less plain and Other examples. The most beautiful object in nature, perhaps, is the countenance of a young and beautiful woman; and we are apt at first to imagine, that, independently of all associations, the forms and colours which it displays are in themselves lovely and engaging, and would appear charming to all beholders, with whatever other qualities or impressions they might happen to be connected. A very little reflection, however, will probably be sufficient


\section*{Source record}
\begin{description}
\item[Scan ID] \texttt{britannica:309f765b1ecc83e0a7f7762e44b4a927a7555443:train:1828}
\item[Image file] \texttt{page-1828-9ca50fbf1eeb.jpg}
\item[Image SHA-256] \texttt{9ca50fbf1eeb13c6abdd7e6d93a69f6c9ad3460d0a6e26c49a0ad5c52b788aa2}
\item[Dataset] NationalLibraryOfScotland/encyclopaedia\_britannica\_illustrated at revision \texttt{309f765b1ecc83e0a7f7762e44b4a927a7555443} (CC0-1.0)
\item[Extraction] model \texttt{nvidia/nemotron-3.5-parse}, raw reply \texttt{extract/raw/page-1828-9ca50fbf1eeb.json}
\item[Raw reply SHA-256] \texttt{edf3f1328430761a8c0509ba114d0f6c3f113743ef72a9bc3da8a489daa87922}
\item[Correction] model \texttt{nvidia/Nemotron-3.5-Super-VL}, block crops, raw replies \texttt{correct/raw/crop\_nothink/page-1828-9ca50fbf1eeb-b*.json}; super\_vl: 7
\item[Running header] BEAUTY. 489
\end{description}
\end{document}
